\chapter{Conclusions and Future Work}
\label{cap:conclusions}

\chapterquote{Any fool can write code that a computer can understand. Good programmers write code that humans can understand.}{Martin Fowler}

\section{Conclusions}

The main objective of this work was to incorporate real-time augmented reality functionality into \textit{OBS Studio}, developing a native filter capable of integrating synthetic content directly into a broadcast video stream. That objective has been met.

The result is a functional plugin that extends \textit{OBS Studio} with four distinct operating modes. In the 3D model mode, the system loads geometry in multiple formats via \textbf{ASSIMP}, estimates the ArUco marker pose in real time using OpenCV's PnP algorithm, and renders the model anchored to the marker with the correct orientation and position, applying the necessary chain of transformations to convert between OpenCV and OBS coordinate conventions. In the countdown mode, that same model acts as an analogue countdown clock whose hands update frame by frame and are periodically synchronised with the official time retrieved from the \textbf{DOMjudge} \citep{DOMjudge} API. In scoreboard mode, the system projects the contest leaderboard onto the video in real time, with all team data organised in four columns. In team info mode, the camera automatically detects which team is in the field of view and displays only that team's data, combining marker detection with a JSON mapping file prepared by the operator before the broadcast.

Beyond the visualisation modes, the work resolved several significant technical problems. The coordinate system gap between OpenCV and OBS required an explicit convention transformation via a $180^\circ$ rotation around the horizontal axis, with the consequent negation of the $Y$ and $Z$ components of the ArUco rotation vector. Overlay smoothing via an exponential moving-average filter — with special handling of angle wrap-around in circular interpolation — eliminates the perceptible jitter caused by camera noise and video compression. The two-thread architecture, with a dedicated secondary thread for HTTP requests protected by a \texttt{CRITICAL\_SECTION}, guarantees that network operations never block rendering. The camera calibration process, implemented as a standalone Python tool, generates the parameter file that the plugin loads on start-up and uses to correct optical distortion for each lens.

The video format conversion system covers the seven formats that OBS can deliver (I420, NV12, YUY2, UYVY, BGRA, RGBA, Y800), making detection work regardless of the connected camera source. The user interface integrated into OBS lets the operator adjust all parameters live, without restarting the scene, which makes the system manageable in a real production environment.

Taken together, the plugin demonstrates that it is possible to add a complete data-driven augmented reality system to \textit{OBS Studio} while maintaining the native filter architecture, with no additional latency from external processes and without modifying the core software.

\section{Limitations}

The plugin was developed and tested on a Windows environment with a single camera. This circumstance implies several limitations worth noting explicitly. The system has not been validated on Linux or macOS, platforms where OBS Studio also runs but whose threading APIs and resource management differ. Detection quality depends directly on lighting conditions and on the print resolution of the markers, factors that have not been evaluated systematically. Camera calibration must be repeated for each different lens, which adds a configuration step that can be impractical in environments where cameras are changed frequently. Finally, the connectivity layer is designed exclusively for the \textbf{DOMjudge} \citep{DOMjudge} REST API, so it cannot connect to other contest management systems without modifying the source code.

\section{Future Work}

The results open several natural lines of continuation.

The most relevant from a technical standpoint is the incorporation of \textbf{markerless tracking}. Current visual SLAM techniques \citep{MurArtal2015} would allow overlays to be anchored to surfaces in the real environment without the need to print and place ArUco codes on the tables, making the system more flexible and applicable to scenarios where the physical space cannot be controlled.

A second line is \textbf{compatibility with other contest management systems}. The current network client is adapted to the DOMjudge API; adding support for platforms such as Codeforces, Kattis, or the ICPC CMS would extend the plugin's reach to a much larger share of the competitive programming community.

On the performance side, it would be valuable to carry out a \textbf{systematic impact evaluation} of the plugin on frame rate and CPU/GPU usage across different hardware configurations, establishing recommended usage thresholds and possible optimisations for mid-range machines.

Finally, adding a \textbf{remote control panel} accessible from a web interface or mobile application would allow the production director to manage modes and overlay visibility from a secondary device, without needing to touch the broadcast machine during the live feed, significantly improving the workflow in productions with more than one person in the technical team.
